%%=============================================================================
%% BIBLIOGRAFIA
%%=============================================================================
\cleardoublepage
\tituloPreliminar{Bibliografía}

La bibliografía es la relación de las fuentes documentales consultadas por el
investigador para sustentar sus trabajos. Su inclusión es obligatoria en todo
trabajo de investigación. Cada referencia bibliográfica se inicia contra el
margen izquierdo.

Existen múltiples formatos de citación y bibliografía por lo cual se está en
libertad de usar cualquier norma bibliográfica de acuerdo con lo acostumbrado
por cada disciplina del conocimiento. En esta medida es necesario que la norma
seleccionada se aplique con rigurosidad. Se recomienda el uso de gestores
bibliográficos como JabRef, KBibTeX, Mendeley o Zotero.

A continuación se lista algunas instituciones que brindan parámetros para el
manejo de las referencias bibliográficas:

{\interlineadoSencillo\small
\begin{longtable}{|p{3.0cm}|p{3.4cm}|p{6.6cm}|}
\hline
\textbf{Institución} & \textbf{Disciplina de aplicación} & \textbf{Vínculos y ejemplos}\\
\hline
\endfirsthead
\hline
\textbf{Institución} & \textbf{Disciplina de aplicación} & \textbf{Vínculos y ejemplos}\\
\hline
\endhead
Modern Language Association (MLA) & Literatura, artes y humanidades. &
  \url{http://www.mla.org/}\newline
  \url{http://biblioteca.udg.es/info_general/Guies/Cites/MLA.asp} (reglamento).\newline
  \url{http://www.liunet.edu/cwis/cwp/library/workshop/citmla.htm} (ejemplos).\\
\hline
American Psychological Association (APA) & Ámbito de la salud (psicología,
  medicina) y en general en todas las ciencias sociales. &
  \url{http://www.apastyle.org/}\newline
  \url{http://biblioteca.udg.es/info_general/Guies/Cites/citar_llibres.asp} (reglamento).\newline
  \url{http://www.liunet.edu/cwis/cwp/library/workshop/citapa.htm} (ejemplos).\\
\hline
Universidad de Chicago/Turabian & Periodismo, historia y humanidades. &
  \url{http://www.chicagomanualofstyle.org/}\newline
  \url{http://www.bedfordstmartins.com/hacker/resdoc/history/footnotes.htm} (Reglamento I)\newline
  \url{http://www.bedfordstmartins.com/online/cite7.html} (Reglamento II).\newline
  \url{http://www.liunet.edu/cwis/cwp/library/workshop/citchi.htm} (ejemplos).\\
\hline
AMA (Asociación Médica de los Estados Unidos) & Ámbito de la salud (psicología,
  medicina) &
  \url{http://healthlinks.washington.edu/hsl/styleguides/ama.htm} (manual de
  estilo de la AMA, que sirve de estándar para las disciplinas que se ocupan de
  medicina, salud y ciencias biológicas).\newline
  \url{http://www.liunet.edu/cwis/cwp/library/workshop/citama.htm} (ejemplos).\\
\hline
Vancouver & Todas las disciplinas &
  \url{http://www.fisterra.com/recursos_web/mbe/vancouver.asp} (estilo de
  Vancouver 2000).\\
\hline
Council of Science Editors (CSE) & En la actualidad abarca diversas ciencias. &
  Scientific Style and Format: The CSE Manual for Authors, Editors, and
  Publishers. 7th ed. 2008. The Ohio State University Libraries. (ejemplos del
  estilo)\\
\hline
National Library of Medicine (NLM) (Biblioteca Nacional de Medicina) & En el
  ámbito médico y, por extensión, en ciencias. &
  \url{http://www.nlm.nih.gov}\newline
  \url{http://www.nlm.nih.gov/pubs/formats/recommendedformats.html} (formatos
  recomendados)\\
\hline
Manual & Todas las disciplinas &
  Patrias, K.: National Library of Medicine recommended formats for
  bibliographic citation. Bethesda (Maryland, EE. UU.): National Library of
  Medicine (Reference Section), 1991.\\
\hline
Pequeño manual sobre fichas bibliográficas & Todas las disciplinas &
  \url{http://www.nuevamuseologia.com.ar/fichasbiblio.htm}\\
\hline
Harvard System of Referencing Guide & Todas las disciplinas &
  Disponibles en: \url{http://libweb.anglia.ac.uk/referencing/harvard.htm}\\
\hline
JabRef y KBibTeX & Todas las disciplinas &
  Herramientas de \LaTeX{} para la gestión de referencias bibliográficas.\\
\hline
\end{longtable}
\par}

\bibliographystyle{unsrt}
\bibliography{referencias}
